\begin{table}[h!]
\centering\small
\caption{Can the magnitude of $\beta_{\mathrm{ME}}$ be reported within the candidate periods?}
\label{tab:period-beta}
\begin{tabular}{p{2.05cm}p{0.75cm}p{2.55cm}p{2.0cm}p{2.0cm}p{3.25cm}}
\toprule
\textbf{Candidate span} & \textbf{$N$} & \textbf{Boundary evidence} & \textbf{$\hat\beta_{\mathrm{ME},j}$} & \textbf{Confidence interval} & \textbf{Empirical verdict}\\
\midrule
1920--1944 & 25 & Initial candidate span & Not estimated & Not available & Magnitude not identified \\
1945--1970 & 26 & Boundary not validated & Not estimated & Not available & Magnitude not identified \\
1971--1988 & 18 & Boundary not validated & Not estimated & Not available & Magnitude not identified \\
1989--2024 & 36 & Boundary not validated & Not estimated & Not available & Magnitude not identified \\
\bottomrule
\multicolumn{6}{p{13.1cm}}{\footnotesize \textbf{Note:} The spans use 1945, 1971, and 1989 only to display the sample sizes implied by the earlier candidate dates. They are not estimated separately. Reporting a coefficient would condition on unvalidated boundaries and mistake persistence for a long-run relation.}
\end{tabular}
\end{table}
